\documentclass[11pt,a4paper]{article}
\usepackage[margin=1in]{geometry}
\usepackage[T1]{fontenc}
\usepackage[utf8]{inputenc}
\usepackage{lmodern}
\usepackage{amsmath,amssymb}
\usepackage{booktabs}
\usepackage{enumitem}
\usepackage{hyperref}
\usepackage{xcolor}
\usepackage{listings}
\usepackage{graphicx}
\graphicspath{{assets/}}

\hypersetup{
  colorlinks=true,
  linkcolor=blue!60!black,
  urlcolor=blue!60!black
}

\lstdefinestyle{pycode}{
  language=Python,
  basicstyle=\ttfamily\small,
  keywordstyle=\color{blue!70!black},
  commentstyle=\color{green!40!black},
  stringstyle=\color{orange!50!black},
  showstringspaces=false,
  breaklines=true,
  frame=single,
  aboveskip=4pt,
  belowskip=4pt
}
\newcommand{\img}[2][]{\includegraphics[#1]{\detokenize{#2}}}
\setlength{\parskip}{2pt}
\setlength{\parindent}{0pt}
\setlength{\abovedisplayskip}{6pt}
\setlength{\belowdisplayskip}{6pt}
\setlength{\abovedisplayshortskip}{4pt}
\setlength{\belowdisplayshortskip}{4pt}
\setlength{\textfloatsep}{8pt}
\renewcommand{\arraystretch}{0.98}
\setlist{nosep}

\begin{document}
\begin{center}
{\LARGE\bfseries Creative Glow \& Tone/Brightness Filters\par}
\vspace{2pt}
{\large xotiena00\par}
\vspace{2pt}
{May 9, 2026\par}
\end{center}
\vspace{6pt}

\section{Implementation and submission}\label{sec:impl}
The full source code is submitted with the project (repository: \texttt{facefiltering}). Filter cores are in \texttt{FaceFiltering/facefiltering/filters/}: \texttt{bloom.py}, \texttt{orton\_effect.py}, \texttt{gamma.py}, \texttt{dodge.py}. Shared utilities: \texttt{facefiltering/ops.py} (Gaussian kernel, convolution), \texttt{facefiltering/validate.py} (BGR \texttt{uint8} checks, clamps), \texttt{facefiltering/registry.py} (\texttt{apply\_filter}). The Gradio UI is \texttt{FaceFiltering/app.py}; batch/CLI use \texttt{FaceFiltering/main.py}. Annex~\ref{sec:annex} lists the correspondence between the equations in this document and these files.

\textbf{Complexity convention.} Costs are stated \emph{per operation} and \emph{for the whole frame}. Let $W\times H$ be the image size, $N_{\mathrm{sp}}=WH$ spatial pixels, $C=3$ BGR channels, and $N_{\mathrm{pix}}=N_{\mathrm{sp}}C$ scalar samples. Point-wise filters (LUT gamma, dodge arithmetic) do $\Theta(1)$ work per scalar sample, hence $\Theta(N_{\mathrm{pix}})$ for the full image. Gaussian blur with a naive $k\times k$ convolution uses $\Theta(k^2)$ multiply-adds per spatial pixel per channel, hence $\Theta(N_{\mathrm{pix}}k^2)$ overall; an optimized separable Gaussian would be $\Theta(N_{\mathrm{pix}}k)$. Saying ``\(O(1)\) per pixel'' for gamma is shorthand for ``per sample''; total work is linear in the number of samples, i.e.\ \(\Theta(N_{\mathrm{pix}})\), not independent of image size.

\section{Classification and Comparison}
\begin{center}
\begin{tabular}{@{}lllll@{}}
\toprule
\textbf{Filter} & \textbf{Group} & \textbf{Methodology} & \textbf{Linear?} & \textbf{Edge Preservation} \\
\midrule
Bloom & Creative Glow & Bright-mask + blur + additive blend & Hybrid & Medium \\
Orton effect & Creative Glow & Blur + screen blend + mix & Hybrid & Medium \\
Gamma & Tone/Brightness & Point-wise power-law mapping (LUT) & No & N/A \\
Dodge & Tone/Brightness & Point-wise nonlinear brightening curve & No & N/A \\
\bottomrule
\end{tabular}
\end{center}

\section{Theory and Application}

\subsection{Bloom}
\textbf{Core model}
\[
M(x)=\mathbf{1}_{I_g(x)\ge T},\quad B=I\odot M,\quad G=G_{\sigma}*B,\quad I'=I+\alpha G
\]
\textbf{Notation} \(T\): threshold, \(\sigma\): blur spread, \(\alpha\): glow strength, \(I'\): output.

\textbf{Pixel-level transformation}
\[
I'_c(x,y)=I_c(x,y)+\alpha\sum_{i=-r}^{r}\sum_{j=-r}^{r}G_{\sigma}(i,j)\,B_c(x+i,y+j)
\]
Only bright pixels feed the glow layer; high \(T\) can suppress bloom even with high \(\alpha\).

\textbf{Algorithm (implementation)}
\begin{lstlisting}[style=pycode]
bright = image * (gray >= T)
glow = convolve(bright, gaussian_kernel(k, sigma))
out = clip(image + alpha * glow)
\end{lstlisting}
\textbf{Implementation detail (steps)} 
(1) Input BGR \(I\) as \texttt{uint8}; compute grayscale \(I_g\) from BGR for the mask only (fixed linear luminance weights matching the code). 
(2) Binary mask \(M=\mathbf{1}_{I_g\ge T}\); bright layer \(B_c=I_c\cdot M\) per channel (float).
(3) Convolution \(G=B*G_\sigma\) with discrete normalized Gaussian stencil of size \(k\times k\), \(\sum G_\sigma=1\).
(4) Output \(I'=\mathrm{clip}_{[0,255]}(I+\alpha G)\) in float then round to \texttt{uint8}.

\textbf{Experiment -- Bloom}
\begin{center}
\begin{tabular}{ccc}
\img[width=0.31\textwidth]{bloom1.jpg} &
\img[width=0.31\textwidth]{bloom2.jpg} &
\img[width=0.31\textwidth]{bloom3.jpg} \\
\scriptsize $T=24,\ \alpha=0.05,\ \sigma=0.5$ &
\scriptsize $T=21,\ \alpha=0.65,\ \sigma=1.2$ &
\scriptsize $T=144,\ \alpha=1.8,\ \sigma=11.3$ \\
\img[width=0.31\textwidth]{bloom1graph.jpg} &
\img[width=0.31\textwidth]{bloom2graph.jpg} &
\img[width=0.31\textwidth]{bloom3graph.jpg}
\end{tabular}
\end{center}

\textbf{Use} cinematic highlight glow and light-source emphasis.  
\begin{center}
\begin{tabular}{ccc}
\img[width=0.31\textwidth]{bloomsample1.jpg} &
\img[width=0.31\textwidth]{bloomsample2.jpg} &
\img[width=0.31\textwidth]{bloomsample3.jpg}
\end{tabular}
\end{center}
\textbf{Compare} stronger spatial glow than gamma/dodge.

\textbf{Cost} Naive Gaussian: \(\Theta(k^2)\) multiply-adds per spatial sample per channel; full image \(\Theta(N_{\mathrm{pix}} k^2)\). Addition \(\Theta(N_{\mathrm{pix}})\). Separable Gaussian (if used) would reduce blur to \(\Theta(N_{\mathrm{pix}} k)\).

\subsection{Orton Effect}
\textbf{Core model}
\[
S = 255 - \frac{(255-I)(255-B)}{255},\quad B=G_{\sigma}*I,\quad I'=(1-\alpha)I+\alpha S
\]
\textbf{Notation} \(B\): blurred image, \(S\): screen blend, \(\sigma\): blur spread, \(\alpha\): mix strength.

\textbf{Pixel-level transformation}
\[
\begin{aligned}
B_c(x,y)&=\sum_{i=-r}^{r}\sum_{j=-r}^{r}G_{\sigma}(i,j)\,I_c(x+i,y+j),\\
S_c(x,y)&=255-\frac{(255-I_c(x,y))(255-B_c(x,y))}{255},\\
I'_c(x,y)&=(1-\alpha)I_c(x,y)+\alpha S_c(x,y).
\end{aligned}
\]

\textbf{Algorithm (implementation)}
\begin{lstlisting}[style=pycode]
blur = convolve(image, gaussian_kernel(k, sigma))
screen = 255 - ((255-image)*(255-blur)/255)
out = clip((1-alpha)*image + alpha*screen)
\end{lstlisting}
\textbf{Implementation detail (steps)} 
(1) Blur full-color image: \(B=I*G_\sigma\) per channel in float. 
(2) Screen blend per channel: \(S_c=255-(255-I_c)(255-B_c)/255\) (8-bit screen in \([0,255]\)). 
(3) Convex mix: \(I'=(1-\alpha)I+\alpha S\), clip and cast to \texttt{uint8}.

\textbf{Experiment -- Orton}
\begin{center}
\begin{tabular}{ccc}
\img[width=0.31\textwidth]{orton1.jpg} &
\img[width=0.31\textwidth]{orton2.jpg} &
\img[width=0.31\textwidth]{orton3.jpg} \\
\scriptsize $\alpha=0.15,\ \sigma=0.1$ &
\scriptsize $\alpha=0.45,\ \sigma=8.5$ &
\scriptsize $\alpha=1.0,\ \sigma=19.1$ \\
\img[width=0.31\textwidth]{orton1graph.jpg} &
\img[width=0.31\textwidth]{orton2graph.jpg} &
\img[width=0.31\textwidth]{orton3graph3.jpg}
\end{tabular}
\end{center}

\textbf{Use} soft-focus portrait look with reduced micro-contrast.  
\begin{center}
\begin{tabular}{ccc}
\img[width=0.31\textwidth]{ortonsample1.jpg} &
\img[width=0.31\textwidth]{ortonsample2.jpg} &
\img[width=0.31\textwidth]{ortonsample3.jpg}
\end{tabular}
\end{center}
\textbf{Compare} gentler and more diffuse than bloom.

\textbf{Cost} Same blur dominance as Bloom: \(\Theta(N_{\mathrm{pix}} k^2)\) for naive \(k\times k\) Gaussian; screen and mix are \(\Theta(N_{\mathrm{pix}})\).

\subsection{Gamma Filter}
\textbf{Core model}
\[
V_{\text{out}} = A\,V_{\text{in}}^{\gamma},\qquad
I' = 255\left(\frac{I}{255}\right)^{1/\gamma}
\]
\textbf{Notation} \(\gamma\): tone-curve exponent, \(A\): scale (usually 1), \(I\in[0,255]\): 8-bit input.
\textbf{Why \(0..255\)} 8-bit \texttt{uint8} images have 256 levels; divide by 255 to normalize to \([0,1]\), then multiply by 255 to return.

\textbf{Pixel-level transformation}
\[
I'_{c}(x,y)=255\left(\frac{I_{c}(x,y)}{255}\right)^{1/\gamma}
\]
\(\gamma<1\): brighter midtones, \(\gamma>1\): darker midtones.

\textbf{Algorithm (implementation)}
\begin{lstlisting}[style=pycode]
lut[i] = round(((i/255.0)**(1.0/gamma))*255)
out = lut[image]
\end{lstlisting}
\textbf{Implementation detail (steps)} 
(1) Clamp \(\gamma\) to a safe interval; precompute LUT \texttt{lut} of length 256: \(\texttt{lut}[i]=\mathrm{round}(255\cdot(i/255)^{1/\gamma})\). 
(2) Apply \texttt{lut} by indexing: each channel \(I'_c(x,y)=\texttt{lut}[I_c(x,y)]\) (vectorized as \texttt{lut[bgr]} in code). 
(3) Output remains BGR \texttt{uint8}.

\textbf{Colour space} The same power law is applied \textbf{independently to each RGB/BGR channel}. This is simple and fast but can shift apparent hue/saturation because chroma is not held fixed. A common alternative is to convert to a luminance-based representation (e.g.\ linear RGB \(\to\) \(Y\)), apply the tone curve to \(Y\) only, and map back, which better preserves perceived color.

\textbf{Experiment -- Gamma}
\begin{center}
\begin{tabular}{ccc}
\img[width=0.31\textwidth]{gamma1.jpg} &
\img[width=0.31\textwidth]{gamma2.jpg} &
\img[width=0.31\textwidth]{gamma3.jpg} \\
\scriptsize $\gamma=0.2$ & \scriptsize $\gamma=1.0$ & \scriptsize $\gamma=1.95$ \\
\img[width=0.31\textwidth]{gamma1graph.jpg} &
\img[width=0.31\textwidth]{gamma2graph.jpg} &
\img[width=0.31\textwidth]{gamma3graph.jpg}
\end{tabular}
\end{center}

\textbf{Use} controlled global tone and midtone correction.  
\begin{center}
\begin{tabular}{ccc}
\img[width=0.31\textwidth]{gammasample1.jpg} &
\img[width=0.31\textwidth]{gammasample2.jpg} &
\img[width=0.31\textwidth]{gammasample3.jpg}
\end{tabular}
\end{center}
\textbf{Compare} smoother and less aggressive than dodge.

\textbf{Cost} \(\Theta(1)\) per scalar lookup; full image \(\Theta(N_{\mathrm{pix}})\). Building the LUT is \(\Theta(1)\) (256 entries).

\subsection{Dodge Filter}
\textbf{Core model}
\[
I' \approx \frac{255\,I}{255-sI}
\]
\textbf{Notation} \(s\): dodge strength, \(I\): input intensity, \(I'\): output intensity.

\textbf{Pixel-level transformation}
\[
I'_{c}(x,y)=\mathrm{clip}\!\left(\frac{255\,I_{c}(x,y)}{\max(255-s\,I_{c}(x,y),1)}\right)
\]
Higher \(s\) strongly boosts highlights and may clip to white.

\textbf{Algorithm (implementation)}
\begin{lstlisting}[style=pycode]
denom = 255 - strength*image
out = clip((image*255) / max(denom, 1))
\end{lstlisting}
\textbf{Implementation detail (steps)} 
(1) Clamp strength \(s\) to a safe range; work in float \([0,255]\). 
(2) Denominator \(d_c=\max(255-s\cdot I_c,\,1)\) per channel per sample (floor at 1 avoids division by zero). 
(3) \(I'_c=\mathrm{clip}(255\cdot I_c/d_c)\), then \texttt{uint8}.

\textbf{Experiment -- Dodge}
\begin{center}
\begin{tabular}{ccc}
\img[width=0.31\textwidth]{dodge1.jpg} &
\img[width=0.31\textwidth]{dodge2.jpg} &
\img[width=0.31\textwidth]{dodge3.jpg} \\
\scriptsize $s=0.01$ & \scriptsize $s=0.5$ & \scriptsize $s=0.90$ \\
\img[width=0.31\textwidth]{dodge1graph.jpg} &
\img[width=0.31\textwidth]{dodge2graph.jpg} &
\img[width=0.31\textwidth]{dodge3graph.jpg}
\end{tabular}
\end{center}

\textbf{Use} stylized highlight brightening and sketch-like effects.  
\begin{center}
\begin{tabular}{ccc}
\img[width=0.31\textwidth]{dodgesample1.jpg} &
\img[width=0.31\textwidth]{dodgesample2.jpg} &
\img[width=0.31\textwidth]{dodgesample3.jpg}
\end{tabular}
\end{center}
\textbf{Compare} more aggressive than gamma; non-spatial (no blur spread).

\textbf{Cost} \(\Theta(1)\) arithmetic per scalar sample; \(\Theta(N_{\mathrm{pix}})\) for the frame.

\subsection{Creative Glow Comparison}
\begin{center}
\begin{tabular}{@{}llll@{}}
\toprule
\textbf{Filter} & \textbf{Visual Effect} & \textbf{Best Use} & \textbf{Cost} \\
\midrule
Bloom & Strong highlight halo/spread & Dramatic bright sources & Medium \\
Orton effect & Soft dreamy glow + lower micro-contrast & Portrait stylization & Medium \\
\bottomrule
\end{tabular}
\end{center}

\subsection{Tone / Brightness Comparison}
\begin{center}
\begin{tabular}{@{}llll@{}}
\toprule
\textbf{Filter} & \textbf{Visual Effect} & \textbf{Best Use} & \textbf{Risk} \\
\midrule
Gamma & Controlled global tone remapping & Midtone correction & Can flatten contrast if overused \\
Dodge & Aggressive highlight boost & Stylized brightening & Highlight clipping/saturation \\
\bottomrule
\end{tabular}
\end{center}

\section{Annex: equation--code map}\label{sec:annex}
\begin{tabular}{@{}ll@{}}
\toprule
\textbf{Concept in this PDF} & \textbf{Implementation file / symbol} \\
\midrule
Bloom (\(M,B,G,G_\sigma,\alpha,\mathrm{clip}\)) & \texttt{facefiltering/filters/bloom.py}, \texttt{apply(...)} \\
Orton blur + screen + mix & \texttt{facefiltering/filters/orton\_effect.py}, \texttt{apply(...)} \\
Gamma LUT \(i\mapsto \mathrm{round}(255(i/255)^{1/\gamma})\) & \texttt{facefiltering/filters/gamma.py}, \texttt{apply(...)} \\
Dodge \(\mathrm{clip}(255I/\max(255-sI,1))\) & \texttt{facefiltering/filters/dodge.py}, \texttt{apply(...)} \\
Gaussian stencil + convolution & \texttt{facefiltering/ops.py} \\
Dispatch \texttt{apply\_filter(name, bgr, **kw)} & \texttt{facefiltering/registry.py} \\
\bottomrule
\end{tabular}

\section{Conclusion}
The two groups serve different goals. Bloom and Orton provide glow-style stylization. Gamma and Dodge remain point-wise tone
operators: gamma for controlled global tone correction and dodge for stronger highlight emphasis.
Gamma and Dodge scale as \(\Theta(N_{\mathrm{pix}})\) over the image; Bloom and Orton add spatial blur at \(\Theta(N_{\mathrm{pix}}k^2)\) per channel for a naive \(k\times k\) kernel (Section~\ref{sec:impl}). The submitted source files in Annex~\ref{sec:annex} match the equations and listings above.

\end{document}
